%=============================================================================
%  report_ieee14_switching_proof_summary_en.tex
%  Technical proof summary, 1--2 pages: why IEEE-14 does not pin all four GFM.
%  Every number is taken from report_ieee14_switching_decision_basis_en.tex.
%
%  BUILD:  xelatex report_ieee14_switching_proof_summary_en.tex   (twice)
%=============================================================================
\documentclass[11pt,a4paper]{article}
\usepackage[a4paper,top=1.3cm,bottom=1.2cm,left=1.7cm,right=1.7cm]{geometry}
\usepackage{newtxtext,newtxmath}
\usepackage{amsmath,booktabs,array,float}
\usepackage[table]{xcolor}
\usepackage{enumitem,fancyhdr,caption,titlesec,microtype}

\definecolor{navy}{RGB}{25,55,95}
\definecolor{ok}{RGB}{0,110,70}
\definecolor{no}{RGB}{170,40,35}
\definecolor{gray}{RGB}{90,98,108}

\pagestyle{fancy}\fancyhf{}
\fancyhead[L]{\small\color{navy}\textbf{Technical proof summary: why not pin all four GFM}}
\fancyhead[R]{\small\color{gray}\thepage}
\renewcommand{\headrulewidth}{0.5pt}
\setlength{\headheight}{14pt}
\setlength{\parindent}{0pt}
\setlength{\parskip}{3pt}
\setlist[itemize]{leftmargin=1.4em,itemsep=0pt,topsep=2pt}
\setlist[enumerate]{leftmargin=1.6em,itemsep=0pt,topsep=2pt}
\captionsetup{font=small,labelfont=bf,skip=1pt}
\titleformat{\section}{\normalfont\normalsize\bfseries\color{navy}}{\thesection}{0.6em}{}
\titlespacing{\section}{0pt}{6pt}{2pt}
\titleformat{\subsection}{\normalfont\small\bfseries\color{navy}}{\thesubsection}{0.5em}{}
\titlespacing{\subsection}{0pt}{4pt}{1pt}

\newcommand{\pu}{\,\mathrm{pu}}
\newcommand{\code}[1]{\texttt{\small #1}}
\newcommand{\zworst}{\zeta_{\mathrm{worst}}}
\newcommand{\zmin}{\zeta_{\min}}
\newcommand{\mes}{\textcolor{ok}{\textbf{[measured]}}}

\begin{document}

\begin{center}
{\large\bfseries\color{navy} Technical proof summary: why IEEE-14 does not pin all four GFM after the SG trip}\\[2pt]
{\small\color{gray} Evidence from \code{report\_ieee14\_switching\_decision\_basis\_en} --- every number read from a named artefact}
\end{center}

\vspace{-4pt}
\noindent
When the SG at bus 1 trips, four dual-mode GFL/GFM converters remain at buses
2, 3, 6 and 8. Why is the forming set not pinned to all four from start to finish,
but chosen adaptively at one or two? There are two measured proofs.

%=============================================================================
\section{Proof 1: SCR does not respond to the forming count}
%=============================================================================

\noindent
The project defines (\code{+stability/ibr\_scr\_metrics.m})
\begin{equation}
\mathrm{SCR}_i = \frac{|V_i|^2}{|Z_{th,ii}|\,S_{\mathrm{rated},i}},
\qquad
Z_{th,ii} = \bigl[\,\mathbf{Y}^{-1}\mathbf{e}_i\,\bigr]_i ,
\label{eq:scr}
\end{equation}
where $\mathbf{Y}$ is built from branch data (including tap and phase shift) plus
bus shunts, \emph{excluding} load admittance, with reference buses grounded. No
term in $\mathbf{Y}$ depends on how many converters are forming. \mes

\begin{table}[H]
\centering\small
\begin{tabular}{@{}lrrrrr@{}}
\toprule
\textbf{Window} & $n$ & \textbf{min} & \textbf{mean} & \textbf{max} & \textbf{$<3$}\\
\midrule
pre-trip (SG online)   & $168$   & $2.0892$ & $2.3439$ & $2.4803$ & $100.00\%$\\
islanded (SG offline)  & $2760$  & $0.6501$ & $\mathbf{2.3015}$ & $2.8147$ & $\mathbf{100.00\%}$\\
whole run, 4 converters & $14716$ & $0.2944$ & $2.5271$ & $6.6726$ & $\mathbf{98.61\%}$\\
\bottomrule
\end{tabular}
\caption{Per-sample SCR at the converter buses. The islanded window averages
$2.3015$ and sits below the WECC strong-grid floor of $3$ for $100.00\%$ of its
samples ($98.61\%$ over the whole run). The minimum $0.2944$ occurs at the bus-9
fault, not at the machine trip. \mes}
\end{table}

\noindent
The first two rows carry the point: the window mean moves from $2.3439$ to
$2.3015$. The machine trip barely moves SCR, because $Z_{th}$ does not know the
machine exists. What changes is \emph{which sources sit behind the bus}, not the
impedance.

\begin{enumerate}
\item \textbf{Promoting a converter cannot stiffen the network.} The impedance
      term of~\eqref{eq:scr} is built from branch data, bus shunts and
      reference-bus grounding only, so adding a former cannot change $Z_{th}$ at
      all. SCR can move only through the measured numerator $|V_i|^2$, never
      through the count. (This does not claim SCR never moves.)
\item \textbf{The problem changes kind, not degree.} REGC\_A screens a
      current-source plant because a PLL needs a stiff bus to lock to. With all
      four forming there is no PLL left to protect, and the question becomes
      whether four voltage sources can share current on a high-impedance network
      --- a question this network answers badly (Proof 2).
\item \textbf{Scope that must travel with the number.} In
      \code{ibr\_scr\_metrics.m:266--276} the gate is forced to profile
      \code{not\_applicable\_\allowbreak full\_state\_\allowbreak source\_model}
      for \code{eecon49\_dual}, giving \code{eligible\_for\_scr = false} and
      \code{pass = true}. The figures above are therefore a per-sample published
      diagnostic, not an active gate. \mes{} The network-invariance conclusion
      does not depend on this: it is a property of the equation, not of the
      gate's state.
\end{enumerate}

%=============================================================================
\section{Proof 2: damping margin and dynamic failure}
%=============================================================================

\noindent
The project enumerates all $2^4-1 = 15$ forming sets, solves each equilibrium and
assembles its full-KCL spectrum. \textbf{Only 7 of the 15 are admissible}, and
\textbf{all four three-unit sets have no equilibrium at all} (the coupled Newton
does not converge; residuals $3.279$, $3.605$, $3.561$, $3.346\times10^{2}$).
What actually gates is equilibrium existence $+$ device limits $+$ full-KCL SSSA
with a damping-ratio floor $\zmin = 0.05$ on every mode of the decision spectrum,
evaluated at three frozen finite-difference steps. $\gamma_{\mathrm{req}} =
0.1$\,rad/s is the \emph{ordering key} among already-admissible candidates only,
not a gate --- the margin may be negative while a set is admissible.

\begin{table}[H]
\centering\small
\begin{tabular}{@{}lrrl@{}}
\toprule
\textbf{Forming set} & $\omega$ (rad/s) & $\zworst$ & \textbf{Outcome}\\
\midrule
$\{5\}$, $\{4\}$, $\{3\}$, $\{2\}$ (one) & $-0.1554$ to $-0.3280$ & $0.2569$--$0.2676$ & admissible\\
$\{2,4\}$, $\{3,4\}$ (two) & $-0.5617$, $\mathbf{-0.5681}$ & $0.2549$, $0.2560$ & admissible (most damped)\\
$\{2,3,4\}$, $\{2,3,5\}$, $\{2,4,5\}$, $\{3,4,5\}$ & --- & --- & \textcolor{no}{\textbf{no equilibrium}}\\
$\{2,3,4,5\}$ (four) & $-0.4829$ & \textcolor{no}{$\mathbf{0.1223}$} & admissible, least damped\\
\bottomrule
\end{tabular}
\caption{$\zworst$ for the other six admissible sets (all four singletons, plus
$\{2,4\}$ and $\{3,4\}$) spans $0.2549$--$0.2676$; the seventh, the four-unit
set, sits at $\mathbf{0.1223}$ --- \textbf{less than half} --- even though it
still clears the $0.05$ floor. \mes}
\end{table}

\noindent
The $0.05$ floor is a minimum, not a target. The four-unit set \emph{passes} the
gate yet is the worst of the six survivors, and going from two units to three is
not a compromise --- it lands where the equations have no solution at all.

\begin{table}[H]
\centering\small
\begin{tabular}{@{}llrrl@{}}
\toprule
\textbf{Arm} & \textbf{horizon} & \textbf{samples} & \textbf{reclose} & \textbf{termination}\\
\midrule
\code{locked\_gfl}  & $20.0000$\,s & $43$   & ---        & \code{noVoltageFormingSource}\\
\code{pinned\_gfm4} & \textcolor{no}{$\mathbf{25.4852}$\,\textbf{s}} & $837$ & never requested & \code{adaptiveDtMin}\\
\code{pinned\_gfm2} & $250.0000$\,s & $1761$ & $154.0244$ & \code{SUCCESS}\\
\code{adaptive}     & $250.0000$\,s & $3679$ & $159.2519$ & \code{SUCCESS}\\
\bottomrule
\end{tabular}
\caption{Six arms from one shared option set. \code{pinned\_gfm4} dies at
$25.4852$\,s --- \textbf{it survives $5.4852$\,s after the SG trip}, before the
load step and before any fault. \mes{} The table also shows that \emph{pinning}
is not the defect: \code{pinned\_gfm2} pins two and reaches $250$\,s. The
evidence condemns the count of four, not the pinned mode.}
\end{table}

\noindent
The terminal state of \code{pinned\_gfm4} is not a gradual loss of damping but
\textbf{frequency splitting}: in its last accepted sample IBR2 runs at
$64.9947$\,Hz while IBR3, IBR6 and IBR8 run at $59.1383$, $59.1809$ and
$59.1946$\,Hz --- \textbf{a $5.86$\,Hz split}, one converter having run away from
the other three. The step controller then reported \textbf{398 successive Newton
non-convergences} and terminated at $\Delta t_{\min} = 6.104\times10^{-6}$\,s.
\mes

\textbf{Why an admissible set can still die:} the acceptance gate is a statement
about an \emph{equilibrium}. It is proved by linearising about the solved
operating point and says nothing about the trajectory that arrives there. The
defect record measured that gap by re-running the same inputs from two different
initial device states.

\begin{table}[H]
\centering\small
\begin{tabular}{@{}llrl@{}}
\toprule
\textbf{Arm} & \textbf{initial device state} & \textbf{outcome} & \textbf{excursion}\\
\midrule
ARM A & the state the production run actually arrives with & \textcolor{no}{\textbf{SLIPS}} & $4358^{\circ}$\\
ARM B & the certified equilibrium state of the same configuration & \textbf{SYNC} & $6.6^{\circ}$\\
\bottomrule
\end{tabular}
\caption{The same configuration, the same inputs, two initial states. The
carrying coordinate is the outer voltage-loop integrator
$gfm\_xi\_Vd/xi\_Vq$ of IBR2: $0.144253$ on arrival against $0.199658$ at the
certified equilibrium --- a $\mathbf{28\%}$ displacement in a coordinate no bus
quantity reflects. \mes}
\end{table}

\noindent
\textbf{Scope of that table.} The ARM A/B experiment is measured on a different
all-four transition from the one that kills \code{pinned\_gfm4}: it is the
load-step commit at $t = 53.4025$\,s in the pre-certificate baseline, not the
death at $25.4852$\,s. Both are the same \emph{class} of failure --- an admitting
equilibrium whose arriving state lies outside its basin of attraction --- but they
are not the same event, and the record does not root-cause the pinned arm's death.
Structurally: the three-unit sets fail the \emph{weaker} test (no equilibrium
exists), while the four-unit set passes the equilibrium test and fails the test of
the \emph{trajectory} --- a failure no eigenvalue check can see.

%=============================================================================
\section{Executive conclusion}
%=============================================================================

\begin{itemize}
\item \textbf{``More GFM is NOT safer'' --- in both directions.} The lower bound
      is measured and hard: \code{locked\_gfl}, with no former at all, dies at
      $t = 20.0000$\,s, the trip instant itself. The upper bound is measured too,
      and it is not four: \code{pinned\_gfm4} dies at $25.4852$\,s.
\item \textbf{Adding GFM does not stiffen the grid.} The impedance term of the
      SCR equation is built from branch data, shunts and reference-bus grounding
      only; the count cannot change it. The islanded window averages $2.3015$ and sits
      below the threshold of $3$ for $100.00\%$ of the islanded time ($98.61\%$
      over the whole run). Pinning all four therefore does not fix the weak grid;
      it moves the problem to four voltage sources sharing current with each
      other.
\item \textbf{The right count is the smallest one still admissible.} Two-unit
      sets are \emph{better damped} than the four-unit set by nearly a factor of
      two ($\zworst$ $0.2549$--$0.2676$ against $0.1223$), and every three-unit
      set has no equilibrium at all. The only exit consistent with the evidence is
      to start at one and grow to two when it is genuinely required --- which is
      what a \emph{minimum-cardinality admissible set} does.
\item \textbf{The counter-evidence belongs in the same answer.} The defect record
      contains a measured case where the four-unit configuration rides a
      $2.0\pu$ load deficit that the two-unit configuration cannot (the defect
      record's load-severity sweep, at a different load level --- a separate
      measurement from the delivered run and from the frozen base-load table).
      The defensible
      conclusion is therefore not ``fewer is better'' but: \textbf{the count is
      not a free parameter and stability is not monotone in it}. It has to be
      chosen per configuration, from the certified margin and from the arrival.
\end{itemize}

\end{document}
